% figs.tex —— 由 ch10/make_figs.py 產生，請勿手動修改。
% 重新產生：python ch10/make_figs.py（第 10 章，動畫 W1001–W1005）

\newcommand{\TenFigTorus}{%
\begin{tikzpicture}[font=\scriptsize,baseline=(current bounding box.center)]
\draw[gray!70,postaction={decorate},decoration={markings,mark=at position 0.56 with {\arrow{Stealth[length=2.6pt]}}}] (-0.190,-0.190) -- (-0.190,1.330);
\draw[gray!70,postaction={decorate},decoration={markings,mark=at position 0.56 with {\arrow{Stealth[length=2.6pt]}}}] (1.330,-0.190) -- (1.330,1.330);
\draw[gray!70,postaction={decorate},decoration={markings,mark=at position 0.44 with {\arrow{Stealth[length=2.6pt]}},mark=at position 0.62 with {\arrow{Stealth[length=2.6pt]}}}] (-0.190,-0.190) -- (1.330,-0.190);
\draw[gray!70,postaction={decorate},decoration={markings,mark=at position 0.44 with {\arrow{Stealth[length=2.6pt]}},mark=at position 0.62 with {\arrow{Stealth[length=2.6pt]}}}] (-0.190,1.330) -- (1.330,1.330);
\fill[gray!55] (0.000,0.000) circle (0.7pt);
\fill[gray!55] (0.380,0.000) circle (0.7pt);
\fill[gray!55] (0.760,0.000) circle (0.7pt);
\fill[gray!55] (1.140,0.000) circle (0.7pt);
\fill[gray!55] (0.000,0.380) circle (0.7pt);
\fill[gray!55] (0.380,0.380) circle (0.7pt);
\fill[gray!55] (0.760,0.380) circle (0.7pt);
\fill[gray!55] (1.140,0.380) circle (0.7pt);
\fill[gray!55] (0.000,0.760) circle (0.7pt);
\fill[gray!55] (0.380,0.760) circle (0.7pt);
\fill[gray!55] (0.760,0.760) circle (0.7pt);
\fill[gray!55] (1.140,0.760) circle (0.7pt);
\fill[gray!55] (0.000,1.140) circle (0.7pt);
\fill[gray!55] (0.380,1.140) circle (0.7pt);
\fill[gray!55] (0.760,1.140) circle (0.7pt);
\fill[gray!55] (1.140,1.140) circle (0.7pt);
\draw[cWalk,very thick,line cap=round] (0.380,0.380) -- (0.000,0.380);
\draw[cWalk,very thick,line cap=round] (0.000,0.380) -- (0.000,0.760);
\draw[cWalk,very thick,line cap=round] (0.000,0.760) -- (0.380,0.760);
\fill[cExact] (0.380,0.380) circle (1.7pt);
\draw[cWalk!70!black,fill=white,thick] (0.380,0.760) circle (1.5pt);
\node[below,align=center,inner sep=1pt] at (0.570,-0.703) {$N=3$：\texttt{WNE}\\不越界};
\draw[gray!70,postaction={decorate},decoration={markings,mark=at position 0.56 with {\arrow{Stealth[length=2.6pt]}}}] (2.660,-0.190) -- (2.660,1.330);
\draw[gray!70,postaction={decorate},decoration={markings,mark=at position 0.56 with {\arrow{Stealth[length=2.6pt]}}}] (4.180,-0.190) -- (4.180,1.330);
\draw[gray!70,postaction={decorate},decoration={markings,mark=at position 0.44 with {\arrow{Stealth[length=2.6pt]}},mark=at position 0.62 with {\arrow{Stealth[length=2.6pt]}}}] (2.660,-0.190) -- (4.180,-0.190);
\draw[gray!70,postaction={decorate},decoration={markings,mark=at position 0.44 with {\arrow{Stealth[length=2.6pt]}},mark=at position 0.62 with {\arrow{Stealth[length=2.6pt]}}}] (2.660,1.330) -- (4.180,1.330);
\fill[gray!55] (2.850,0.000) circle (0.7pt);
\fill[gray!55] (3.230,0.000) circle (0.7pt);
\fill[gray!55] (3.610,0.000) circle (0.7pt);
\fill[gray!55] (3.990,0.000) circle (0.7pt);
\fill[gray!55] (2.850,0.380) circle (0.7pt);
\fill[gray!55] (3.230,0.380) circle (0.7pt);
\fill[gray!55] (3.610,0.380) circle (0.7pt);
\fill[gray!55] (3.990,0.380) circle (0.7pt);
\fill[gray!55] (2.850,0.760) circle (0.7pt);
\fill[gray!55] (3.230,0.760) circle (0.7pt);
\fill[gray!55] (3.610,0.760) circle (0.7pt);
\fill[gray!55] (3.990,0.760) circle (0.7pt);
\fill[gray!55] (2.850,1.140) circle (0.7pt);
\fill[gray!55] (3.230,1.140) circle (0.7pt);
\fill[gray!55] (3.610,1.140) circle (0.7pt);
\fill[gray!55] (3.990,1.140) circle (0.7pt);
\draw[cWalk,very thick,line cap=round] (3.230,0.380) -- (2.850,0.380);
\draw[cWalk,very thick,line cap=round] (2.850,0.380) -- (2.850,0.000);
\draw[cBar,very thick,line cap=round] (2.850,0.000) -- (2.850,-0.190);
\draw[cBar,very thick,line cap=round] (2.850,1.330) -- (2.850,1.140);
\draw[cWalk,very thick,line cap=round] (2.850,1.140) -- (3.230,1.140);
\draw[cWalk,very thick,line cap=round] (3.230,1.140) -- (3.610,1.140);
\draw[cBar,very thick,line cap=round] (3.610,1.140) -- (3.610,1.330);
\draw[cBar,very thick,line cap=round] (3.610,-0.190) -- (3.610,0.000);
\draw[cWalk,very thick,line cap=round] (3.610,0.000) -- (3.610,0.380);
\draw[cWalk,very thick,line cap=round] (3.610,0.380) -- (3.990,0.380);
\fill[cExact] (3.230,0.380) circle (1.7pt);
\draw[cWalk!70!black,fill=white,thick] (3.990,0.380) circle (1.5pt);
\node[below,align=center,inner sep=1pt] at (3.420,-0.703) {$N=8$：\texttt{WSSEENNE}\\越界 2 次};
\draw[gray!70,postaction={decorate},decoration={markings,mark=at position 0.56 with {\arrow{Stealth[length=2.6pt]}}}] (5.510,-0.190) -- (5.510,1.330);
\draw[gray!70,postaction={decorate},decoration={markings,mark=at position 0.56 with {\arrow{Stealth[length=2.6pt]}}}] (7.030,-0.190) -- (7.030,1.330);
\draw[gray!70,postaction={decorate},decoration={markings,mark=at position 0.44 with {\arrow{Stealth[length=2.6pt]}},mark=at position 0.62 with {\arrow{Stealth[length=2.6pt]}}}] (5.510,-0.190) -- (7.030,-0.190);
\draw[gray!70,postaction={decorate},decoration={markings,mark=at position 0.44 with {\arrow{Stealth[length=2.6pt]}},mark=at position 0.62 with {\arrow{Stealth[length=2.6pt]}}}] (5.510,1.330) -- (7.030,1.330);
\fill[gray!55] (5.700,0.000) circle (0.7pt);
\fill[gray!55] (6.080,0.000) circle (0.7pt);
\fill[gray!55] (6.460,0.000) circle (0.7pt);
\fill[gray!55] (6.840,0.000) circle (0.7pt);
\fill[gray!55] (5.700,0.380) circle (0.7pt);
\fill[gray!55] (6.080,0.380) circle (0.7pt);
\fill[gray!55] (6.460,0.380) circle (0.7pt);
\fill[gray!55] (6.840,0.380) circle (0.7pt);
\fill[gray!55] (5.700,0.760) circle (0.7pt);
\fill[gray!55] (6.080,0.760) circle (0.7pt);
\fill[gray!55] (6.460,0.760) circle (0.7pt);
\fill[gray!55] (6.840,0.760) circle (0.7pt);
\fill[gray!55] (5.700,1.140) circle (0.7pt);
\fill[gray!55] (6.080,1.140) circle (0.7pt);
\fill[gray!55] (6.460,1.140) circle (0.7pt);
\fill[gray!55] (6.840,1.140) circle (0.7pt);
\draw[cWalk,very thick,line cap=round] (6.080,0.380) -- (5.700,0.380);
\draw[cWalk,very thick,line cap=round] (5.700,0.380) -- (5.700,0.000);
\draw[cBar,very thick,line cap=round] (5.700,0.000) -- (5.700,-0.190);
\draw[cBar,very thick,line cap=round] (5.700,1.330) -- (5.700,1.140);
\draw[cWalk,very thick,line cap=round] (5.700,1.140) -- (6.080,1.140);
\draw[cWalk,very thick,line cap=round] (6.080,1.140) -- (6.460,1.140);
\draw[cWalk,very thick,line cap=round] (6.460,1.140) -- (6.840,1.140);
\draw[cBar,very thick,line cap=round] (6.840,1.140) -- (6.840,1.330);
\draw[cBar,very thick,line cap=round] (6.840,-0.190) -- (6.840,0.000);
\draw[cWalk,very thick,line cap=round] (6.840,0.000) -- (6.840,0.380);
\draw[cWalk,very thick,line cap=round] (6.840,0.380) -- (6.840,0.760);
\draw[cBar,very thick,line cap=round] (6.840,0.760) -- (7.030,0.760);
\draw[cBar,very thick,line cap=round] (5.510,0.760) -- (5.700,0.760);
\draw[cWalk,very thick,line cap=round] (5.700,0.760) -- (6.080,0.760);
\draw[cWalk,very thick,line cap=round] (6.080,0.760) -- (6.460,0.760);
\draw[cWalk,very thick,line cap=round] (6.460,0.760) -- (6.460,0.380);
\draw[cWalk,very thick,line cap=round] (6.460,0.380) -- (6.460,0.000);
\draw[cWalk,very thick,line cap=round] (6.460,0.000) -- (6.080,0.000);
\fill[cExact] (6.080,0.380) circle (1.7pt);
\draw[cWalk!70!black,fill=white,thick] (6.080,0.000) circle (1.5pt);
\node[below,align=center,inner sep=1pt] at (6.270,-0.703) {$N=15$：走滿 16 格\\越界 3 次};
\begin{scope}
\clip (8.415,-0.665) rectangle (10.125,1.425);
\draw[gray!60,densely dashed] (8.510,-0.190) rectangle (10.030,1.330);
\draw[gray!60,densely dashed] (8.510,-1.710) rectangle (10.030,-0.190);
\fill[gray!40] (8.700,-0.380) circle (0.7pt);
\fill[gray!40] (9.080,-0.380) circle (0.7pt);
\fill[gray!40] (9.460,-0.380) circle (0.7pt);
\fill[gray!40] (9.840,-0.380) circle (0.7pt);
\fill[gray!40] (8.700,0.000) circle (0.7pt);
\fill[gray!40] (9.080,0.000) circle (0.7pt);
\fill[gray!40] (9.460,0.000) circle (0.7pt);
\fill[gray!40] (9.840,0.000) circle (0.7pt);
\fill[gray!40] (8.700,0.380) circle (0.7pt);
\fill[gray!40] (9.080,0.380) circle (0.7pt);
\fill[gray!40] (9.460,0.380) circle (0.7pt);
\fill[gray!40] (9.840,0.380) circle (0.7pt);
\fill[gray!40] (8.700,0.760) circle (0.7pt);
\fill[gray!40] (9.080,0.760) circle (0.7pt);
\fill[gray!40] (9.460,0.760) circle (0.7pt);
\fill[gray!40] (9.840,0.760) circle (0.7pt);
\fill[gray!40] (8.700,1.140) circle (0.7pt);
\fill[gray!40] (9.080,1.140) circle (0.7pt);
\fill[gray!40] (9.460,1.140) circle (0.7pt);
\fill[gray!40] (9.840,1.140) circle (0.7pt);
\end{scope}
\draw[cWalk,very thick,line join=round,line cap=round] (9.080,0.380) -- (8.700,0.380) -- (8.700,0.000) -- (8.700,-0.380) -- (9.080,-0.380) -- (9.460,-0.380) -- (9.460,0.000) -- (9.460,0.380) -- (9.840,0.380);
\fill[cExact] (9.080,0.380) circle (1.7pt);
\draw[cWalk!70!black,fill=white,thick] (9.840,0.380) circle (1.5pt);
\node[below,align=center,inner sep=1pt] at (9.270,-0.703) {$N=8$ 攤開到平面\\終點 $(3,1)$};
\end{tikzpicture}}

\newcommand{\TenTabTorus}{%
\begin{tabular}{@{}l*{8}{r}@{}}
\toprule
$N$ & 0 & 1 & 2 & 3 & 4 & 5 & 6 & 7\\
\midrule
$4\times4$ 環面 $c_N$ & {\color{cExact}1} & {\color{cExact}4} & {\color{cExact}12} & {\color{cExact}36} & \textbf{96} & 264 & 624 & 1536\\
無限平面 $c_N$ & {\color{cExact}1} & {\color{cExact}4} & {\color{cExact}12} & {\color{cExact}36} & \textbf{100} & 284 & 780 & 2172\\
\midrule
$N$ & 8 & 9 & 10 & 11 & 12 & 13 & 14 & 15\\
\midrule
$4\times4$ 環面 $c_N$ & 3072 & 6240 & 9840 & 15912 & 17472 & 19200 & 10656 & 5712\\
無限平面 $c_N$ & 5916 & 16268 & 44100 & 120292 & 324932 & 881500 & 2374444 & 6416596\\
\bottomrule
\end{tabular}}

\newcommand{\TenFigMF}{%
\begin{tikzpicture}[x=5.0cm,y=1.9cm,>=Stealth,font=\scriptsize,baseline=(current bounding box.center)]
\draw[->] (0,0) -- (1.12,0) node[right] {$\lambda$};
\draw[->] (0,0) -- (0,1.2) node[above] {$f(\lambda)/$峰值};
\foreach \x in {0,0.5,1} {\draw (\x,0) -- (\x,-0.03) node[below] {$\x$};}
\foreach \y in {0.5,1} {\draw (0,\y) -- (-0.015,\y) node[left] {$\y$};}
\draw[gray!60,densely dotted] (0.5,0) -- (0.5,1.06);
\draw[cBar,thick] plot coordinates {(0.0000,1.0000) (0.0067,0.8455) (0.0133,0.7134) (0.0200,0.6005) (0.0267,0.5044) (0.0333,0.4227) (0.0400,0.3535) (0.0467,0.2949) (0.0533,0.2455) (0.0600,0.2039) (0.0667,0.1690) (0.0733,0.1398) (0.0800,0.1153) (0.0867,0.0949) (0.0933,0.0780) (0.1000,0.0639) (0.1067,0.0523) (0.1133,0.0427) (0.1200,0.0347) (0.1267,0.0282) (0.1333,0.0229) (0.1400,0.0185) (0.1467,0.0149) (0.1533,0.0120) (0.1600,0.0097) (0.1667,0.0077) (0.1733,0.0062) (0.1800,0.0049) (0.1867,0.0039) (0.1933,0.0031) (0.2000,0.0025) (0.2067,0.0020) (0.2133,0.0015) (0.2200,0.0012) (0.2267,0.0010) (0.2333,0.0008) (0.2400,0.0006) (0.2467,0.0005) (0.2533,0.0004) (0.2600,0.0003) (0.2667,0.0002) (0.2733,0.0002) (0.2800,0.0001) (0.2867,0.0001) (0.2933,0.0001) (0.3000,0.0001) (0.4500,0.0000) (0.6000,0.0000) (0.8000,0.0000) (1.0000,0.0000)};
\draw[cWalk,thick] plot coordinates {(0.0000,0.0019) (0.0169,0.0029) (0.0339,0.0044) (0.0508,0.0065) (0.0678,0.0094) (0.0847,0.0134) (0.1017,0.0189) (0.1186,0.0264) (0.1356,0.0362) (0.1525,0.0489) (0.1695,0.0652) (0.1864,0.0856) (0.2034,0.1109) (0.2203,0.1415) (0.2373,0.1781) (0.2542,0.2209) (0.2712,0.2701) (0.2881,0.3256) (0.3051,0.3868) (0.3220,0.4530) (0.3390,0.5230) (0.3559,0.5952) (0.3729,0.6677) (0.3898,0.7383) (0.4068,0.8047) (0.4237,0.8646) (0.4407,0.9158) (0.4576,0.9561) (0.4746,0.9840) (0.4915,0.9982) (0.5085,0.9982) (0.5254,0.9840) (0.5424,0.9561) (0.5593,0.9158) (0.5763,0.8646) (0.5932,0.8047) (0.6102,0.7383) (0.6271,0.6677) (0.6441,0.5952) (0.6610,0.5230) (0.6780,0.4530) (0.6949,0.3868) (0.7119,0.3256) (0.7288,0.2701) (0.7458,0.2209) (0.7627,0.1781) (0.7797,0.1415) (0.7966,0.1109) (0.8136,0.0856) (0.8305,0.0652) (0.8475,0.0489) (0.8644,0.0362) (0.8814,0.0264) (0.8983,0.0189) (0.9153,0.0134) (0.9322,0.0094) (0.9492,0.0065) (0.9661,0.0044) (0.9831,0.0029) (1.0000,0.0019)};
\node[cBar,anchor=west,align=left] at (0.035,0.66) {$r=0.5$\\（$\vec M\approx0$）};
\node[cWalk,anchor=west,align=left] at (0.66,0.8) {$r=-0.5$：峰在\\$\lambda^*=|r|/u$};
\end{tikzpicture}}

\newcommand{\TenTabN}{%
\begin{tabular}{@{}l*{6}{c}@{}}
\toprule
$r=0.5$，$n=$ & 3 & 2 & 1 & 0.1 & 0.01 & $\to0$\\
\midrule
$\Gamma_n$ & 1.7183 & 1.7634 & 1.8131 & 1.8624 & 1.8676 & 1.8682\\
\midrule
$r=-0.5$，$n=$ & $10^{-2}$ & $10^{-3}$ & $10^{-4}$ & $10^{-6}$ & & \\
\midrule
$\Gamma_n/\Gamma_{MF}$ & 0.3299 & 0.8311 & 0.9801 & 0.9998 & & \\
\bottomrule
\end{tabular}}

\newcommand{\TenTabAsym}{%
\begin{tabular}{@{}lcccc@{}}
\toprule
$N_0$ & 100 & 1000 & $10^4$ & 極限\\
\midrule
$r=0.5$：$\Gamma_{MF}$ & 1.8682 & 1.9844 & 1.9984 & $\frac1r=2$\\
$r=-0.5$：精確 $\div$ 漸近 & 0.9998 & 1.0000 & 1.0000 & 1\\
$r=-0.5$：$\ln\Gamma_{MF}/N_0$ & 0.0912 & 0.0665 & 0.0630 & $\frac{r^2}{4u}=0.0625$\\
\bottomrule
\end{tabular}}

\newcommand{\TenFigWall}{%
\begin{tikzpicture}[x=0.8cm,y=1.9cm,>=Stealth,font=\scriptsize,baseline=(current bounding box.center)]
\fill[gray!22] (-0.32,0) rectangle (0,1.18);
\draw[gray!70!black,thick] (0,0) -- (0,1.18);
\node[gray!60!black,font=\tiny] at (-0.16,0.12) {牆};
\draw[->] (0,0) -- (6.5,0) node[right] {$x$};
\foreach \x in {2,4,6} {\draw (\x,0) -- (\x,-0.03) node[below] {$\x$};}
\foreach \y in {0.5,1} {\draw (-0.32,\y) -- (-0.4,\y) node[left] {$\y$};}
\node[left] at (-0.32,1.12) {$M$};
\draw[cExact,densely dashed] (0,1) -- (6.2,1) node[above left] {$M_0=\sqrt{|r|/u}=1$};
\draw[gray!60,densely dotted] (0,0) -- (1.4142,1) -- (1.4142,0);
\node[anchor=west,gray!60!black] at (1.4142,0.2) {癒合長度 $\sqrt{2/|r|}=1.4142$};
\draw[cBar,thick] plot coordinates {(0.000,0.0000) (0.102,0.0718) (0.203,0.1428) (0.305,0.2124) (0.407,0.2800) (0.508,0.3448) (0.610,0.4065) (0.712,0.4648) (0.814,0.5192) (0.915,0.5698) (1.017,0.6163) (1.119,0.6590) (1.220,0.6978) (1.322,0.7328) (1.424,0.7644) (1.525,0.7927) (1.627,0.8179) (1.729,0.8404) (1.831,0.8603) (1.932,0.8778) (2.034,0.8933) (2.136,0.9070) (2.237,0.9189) (2.339,0.9294) (2.441,0.9386) (2.542,0.9466) (2.644,0.9536) (2.746,0.9597) (2.847,0.9650) (2.949,0.9696) (3.051,0.9736) (3.153,0.9771) (3.254,0.9801) (3.356,0.9828) (3.458,0.9851) (3.559,0.9871) (3.661,0.9888) (3.763,0.9903) (3.864,0.9916) (3.966,0.9927) (4.068,0.9937) (4.169,0.9945) (4.271,0.9952) (4.373,0.9959) (4.475,0.9964) (4.576,0.9969) (4.678,0.9973) (4.780,0.9977) (4.881,0.9980) (4.983,0.9983) (5.085,0.9985) (5.186,0.9987) (5.288,0.9989) (5.390,0.9990) (5.492,0.9992) (5.593,0.9993) (5.695,0.9994) (5.797,0.9994) (5.898,0.9995) (6.000,0.9996)};
\fill[cBar] (1,0.6089) circle (1.4pt); \node[cBar,anchor=north west,inner sep=2pt] at (1,0.6089) {0.6089};
\fill[cBar] (2,0.8884) circle (1.4pt); \node[cBar,anchor=north west,inner sep=2pt] at (2,0.8884) {0.8884};
\fill[cBar] (4,0.9930) circle (1.4pt); \node[cBar,anchor=north,inner sep=2pt] at (4,0.9930) {0.9930};
\end{tikzpicture}}

\newcommand{\TenTabModes}{%
\begin{tabular}{@{}lcccc@{}}
\toprule
$n$ & 3 & 2 & 1 & $\to0$\\
\midrule
橫向 $\lambda_1=\frac{1}{2}$ 的個數 & $2$ & $1$ & $0$ & $-1$\\
縱向 $\lambda_2=\frac{3}{2}$ 的個數 & 1 & 1 & 1 & 1\\
$\det X=\lambda_1^{n-1}\lambda_2$ & $\frac{3}{8}$ & $\frac{3}{4}$ & $\frac{3}{2}$ & $3$\\
$(\det X)^{-1/2}$ & 1.6330 & 1.1547 & 0.8165 & 0.5774\\
\bottomrule
\end{tabular}}

\newcommand{\TenFigHole}{%
\begin{tikzpicture}[x=0.165cm,y=24cm,>=Stealth,font=\scriptsize,baseline=(current bounding box.center)]
\draw[->] (0,0) -- (35,0) node[right] {$\Delta$};
\draw[->] (0,-0.122) -- (0,0.014) node[above] {$-D$};
\foreach \x in {8,16,24,32} {\draw (\x,0) -- (\x,0.002) node[above] {$\x$};}
\node[above left,inner sep=1pt] at (0,0) {$0$};
\foreach \y/\l in {-0.05/-0.05,-0.1/-0.10} {\draw (0,\y) -- (-0.45,\y) node[left] {$\l$};}
\draw[cWalk,densely dashed,thick] plot coordinates {(2,-0.03659) (3,-0.02324) (4,-0.01657) (5,-0.01256) (6,-0.00989) (7,-0.00799) (8,-0.00656) (9,-0.00545) (10,-0.00456) (11,-0.00383) (12,-0.00322) (13,-0.00271) (14,-0.00227) (15,-0.00189) (16,-0.00155) (17,-0.00126) (18,-0.00100) (19,-0.00076) (20,-0.00055) (21,-0.00036) (22,-0.00019) (23,-0.00003) (24,0.00012) (25,0.00025) (26,0.00037) (27,0.00049) (28,0.00059) (29,0.00069) (30,0.00078) (31,0.00087) (32,0.00095)};
\draw[cBar,thick] plot coordinates {(0,-0.10827) (1,-0.06130) (2,-0.03586) (3,-0.02339) (4,-0.01665) (5,-0.01257) (6,-0.00986) (7,-0.00793) (8,-0.00650) (9,-0.00540) (10,-0.00452) (11,-0.00380) (12,-0.00321) (13,-0.00272) (14,-0.00230) (15,-0.00194) (16,-0.00163) (17,-0.00136) (18,-0.00113) (19,-0.00093) (20,-0.00075) (21,-0.00059) (22,-0.00046) (23,-0.00034) (24,-0.00024) (25,-0.00015) (26,-0.00008) (27,-0.00002) (28,0.00003) (29,0.00007) (30,0.00010) (31,0.00011) (32,0.00012)};
\fill[cBar] (0,-0.10827) circle (1.1pt);
\fill[cBar] (1,-0.06130) circle (1.1pt);
\fill[cBar] (2,-0.03586) circle (1.1pt);
\fill[cBar] (4,-0.01665) circle (1.1pt);
\fill[cBar] (8,-0.00650) circle (1.1pt);
\fill[cBar] (16,-0.00163) circle (1.1pt);
\fill[cBar] (32,0.00012) circle (1.1pt);
\node[cBar,anchor=west] at (1.3,-0.1083) {$\Delta=0$：$1-D=0.8917$（少 11\%）};
\draw[cBar,thick] (11,-0.072) -- (14,-0.072) node[right,black] {$-D(\Delta)$，$64^3$ 週期格子};
\draw[cWalk,densely dashed,thick] (11,-0.088) -- (14,-0.088) node[right,black] {$c_0-c/\Delta$，$c=0.0801\approx\frac1{4\pi}$};
\end{tikzpicture}}

\newcommand{\TenTabHole}{%
\begin{tabular}{@{}lcccc@{}}
\toprule
$D$（$r=-1$），$\rho=$ & 1 & 2 & 4 & 8\\
\midrule
連續 $\frac{1-e^{-\sqrt2\rho}}{4\pi\rho}$ & 0.0602 & 0.0374 & 0.0198 & 0.0099\\
無限大格子 & 0.0648 & 0.0394 & 0.0202 & 0.0100\\
$64^3$ 週期格子 & 0.0613 & 0.0359 & 0.0167 & 0.0065\\
Coulomb $\frac1{4\pi\rho}$ & 0.0796 & 0.0398 & 0.0199 & 0.0099\\
\bottomrule
\end{tabular}}

\newcommand{\TenTabVisits}{%
\begin{tabular}{@{}l*{5}{c}@{}}
\toprule
$R$ & 1 & 2 & 4 & 8 & 16\\
\midrule
精確 & 0.4509 & 0.1931 & 0.0610 & 0.0103 & 0.000183\\
erfc 形 & 0.4121 & 0.1744 & 0.0586 & 0.0100 & 0.000174\\
比 & 1.094 & 1.108 & 1.042 & 1.028 & 1.052\\
書上 &  &  &  & {\color{cWrong}$4.3\times10^{-4}$} & {\color{cWrong}$1.7\times10^{-28}$}\\
\midrule
$G(R)\cdot\frac{2\pi R}3$ & 1.0815 & 1.0779 & 1.0198 & 1.0041 & 1.0009\\
\bottomrule
\end{tabular}}

\newcommand{\TenFigRho}{%
\begin{tikzpicture}[x=4.6cm,y=2.5cm,>=Stealth,font=\scriptsize,baseline=(current bounding box.center)]
\draw[->] (0,0) -- (1.29,0) node[right] {$z$};
\draw[->] (0,0) -- (0,1.08);
\foreach \x in {0,0.6,1.2} {\draw (\x,0) -- (\x,-0.025) node[below] {$\x$};}
\foreach \y in {0.5,1} {\draw (0,\y) -- (-0.012,\y) node[left] {$\y$};}
\draw[gray!60,densely dashed] (0.3791,0) -- (0.3791,0.84) node[above,black] {$z_c$};
\draw[cBar!40,thick] plot coordinates {(0.0000,0.0000) (0.0203,0.0094) (0.0407,0.0197) (0.0610,0.0309) (0.0814,0.0431) (0.1017,0.0565) (0.1220,0.0711) (0.1424,0.0870) (0.1627,0.1042) (0.1831,0.1227) (0.2034,0.1424) (0.2237,0.1633) (0.2441,0.1850) (0.2644,0.2075) (0.2847,0.2305) (0.3051,0.2538) (0.3254,0.2772) (0.3458,0.3003) (0.3661,0.3231) (0.3864,0.3454) (0.4068,0.3670) (0.4271,0.3878) (0.4475,0.4077) (0.4678,0.4268) (0.4881,0.4449) (0.5085,0.4620) (0.5288,0.4783) (0.5492,0.4937) (0.5695,0.5082) (0.5898,0.5219) (0.6102,0.5348) (0.6305,0.5469) (0.6508,0.5584) (0.6712,0.5692) (0.6915,0.5794) (0.7119,0.5890) (0.7322,0.5981) (0.7525,0.6067) (0.7729,0.6149) (0.7932,0.6226) (0.8136,0.6299) (0.8339,0.6368) (0.8542,0.6434) (0.8746,0.6497) (0.8949,0.6556) (0.9153,0.6613) (0.9356,0.6667) (0.9559,0.6719) (0.9763,0.6768) (0.9966,0.6815) (1.0169,0.6861) (1.0373,0.6904) (1.0576,0.6945) (1.0780,0.6985) (1.0983,0.7023) (1.1186,0.7060) (1.1390,0.7095) (1.1593,0.7129) (1.1797,0.7161) (1.2000,0.7193)};
\draw[cWalk!40,densely dashed,thick] plot coordinates {(0.0000,0.0000) (0.0203,0.0098) (0.0407,0.0214) (0.0610,0.0351) (0.0814,0.0512) (0.1017,0.0699) (0.1220,0.0915) (0.1424,0.1159) (0.1627,0.1429) (0.1831,0.1721) (0.2034,0.2030) (0.2237,0.2346) (0.2441,0.2660) (0.2644,0.2964) (0.2847,0.3247) (0.3051,0.3503) (0.3254,0.3726) (0.3458,0.3912) (0.3661,0.4059) (0.3864,0.4168) (0.4068,0.4241) (0.4271,0.4280) (0.4475,0.4289) (0.4678,0.4273) (0.4881,0.4235) (0.5085,0.4180) (0.5288,0.4112) (0.5492,0.4032) (0.5695,0.3945) (0.5898,0.3853) (0.6102,0.3757) (0.6305,0.3660) (0.6508,0.3562) (0.6712,0.3465) (0.6915,0.3369) (0.7119,0.3275) (0.7322,0.3183) (0.7525,0.3095) (0.7729,0.3009) (0.7932,0.2926) (0.8136,0.2847) (0.8339,0.2770) (0.8542,0.2696) (0.8746,0.2626) (0.8949,0.2558) (0.9153,0.2493) (0.9356,0.2431) (0.9559,0.2372) (0.9763,0.2314) (0.9966,0.2260) (1.0169,0.2207) (1.0373,0.2157) (1.0576,0.2109) (1.0780,0.2063) (1.0983,0.2018) (1.1186,0.1975) (1.1390,0.1934) (1.1593,0.1895) (1.1797,0.1857) (1.2000,0.1821)};
\fill[cWalk!40] (0.4443,0.4289) circle (1.3pt);
\draw[cBar!70,thick] plot coordinates {(0.0000,0.0000) (0.0203,0.0053) (0.0407,0.0111) (0.0610,0.0176) (0.0814,0.0247) (0.1017,0.0328) (0.1220,0.0419) (0.1424,0.0522) (0.1627,0.0640) (0.1831,0.0774) (0.2034,0.0927) (0.2237,0.1101) (0.2441,0.1296) (0.2644,0.1514) (0.2847,0.1754) (0.3051,0.2013) (0.3254,0.2288) (0.3458,0.2575) (0.3661,0.2868) (0.3864,0.3163) (0.4068,0.3454) (0.4271,0.3737) (0.4475,0.4009) (0.4678,0.4267) (0.4881,0.4511) (0.5085,0.4739) (0.5288,0.4952) (0.5492,0.5150) (0.5695,0.5334) (0.5898,0.5505) (0.6102,0.5663) (0.6305,0.5810) (0.6508,0.5947) (0.6712,0.6075) (0.6915,0.6194) (0.7119,0.6305) (0.7322,0.6409) (0.7525,0.6506) (0.7729,0.6598) (0.7932,0.6684) (0.8136,0.6766) (0.8339,0.6842) (0.8542,0.6915) (0.8746,0.6984) (0.8949,0.7049) (0.9153,0.7111) (0.9356,0.7169) (0.9559,0.7225) (0.9763,0.7279) (0.9966,0.7330) (1.0169,0.7378) (1.0373,0.7425) (1.0576,0.7470) (1.0780,0.7512) (1.0983,0.7553) (1.1186,0.7592) (1.1390,0.7630) (1.1593,0.7666) (1.1797,0.7701) (1.2000,0.7735)};
\draw[cWalk!70,densely dashed,thick] plot coordinates {(0.0000,0.0000) (0.0203,0.0055) (0.0407,0.0122) (0.0610,0.0203) (0.0814,0.0303) (0.1017,0.0427) (0.1220,0.0581) (0.1424,0.0771) (0.1627,0.1005) (0.1831,0.1290) (0.2034,0.1631) (0.2237,0.2029) (0.2441,0.2482) (0.2644,0.2977) (0.2847,0.3497) (0.3051,0.4018) (0.3254,0.4510) (0.3458,0.4948) (0.3661,0.5308) (0.3864,0.5578) (0.4068,0.5752) (0.4271,0.5835) (0.4475,0.5839) (0.4678,0.5776) (0.4881,0.5662) (0.5085,0.5513) (0.5288,0.5339) (0.5492,0.5153) (0.5695,0.4961) (0.5898,0.4769) (0.6102,0.4581) (0.6305,0.4399) (0.6508,0.4226) (0.6712,0.4062) (0.6915,0.3908) (0.7119,0.3762) (0.7322,0.3625) (0.7525,0.3497) (0.7729,0.3377) (0.7932,0.3264) (0.8136,0.3158) (0.8339,0.3058) (0.8542,0.2965) (0.8746,0.2876) (0.8949,0.2793) (0.9153,0.2714) (0.9356,0.2640) (0.9559,0.2570) (0.9763,0.2503) (0.9966,0.2439) (1.0169,0.2379) (1.0373,0.2322) (1.0576,0.2267) (1.0780,0.2215) (1.0983,0.2165) (1.1186,0.2117) (1.1390,0.2072) (1.1593,0.2028) (1.1797,0.1986) (1.2000,0.1946)};
\fill[cWalk!70] (0.4380,0.5846) circle (1.3pt);
\draw[cBar,thick] plot coordinates {(0.0000,0.0000) (0.0203,0.0034) (0.0407,0.0071) (0.0610,0.0112) (0.0814,0.0159) (0.1017,0.0211) (0.1220,0.0270) (0.1424,0.0339) (0.1627,0.0418) (0.1831,0.0512) (0.2034,0.0623) (0.2237,0.0756) (0.2441,0.0913) (0.2644,0.1101) (0.2847,0.1323) (0.3051,0.1580) (0.3254,0.1873) (0.3458,0.2196) (0.3661,0.2542) (0.3864,0.2900) (0.4068,0.3261) (0.4271,0.3612) (0.4475,0.3948) (0.4678,0.4263) (0.4881,0.4554) (0.5085,0.4821) (0.5288,0.5065) (0.5492,0.5288) (0.5695,0.5491) (0.5898,0.5677) (0.6102,0.5847) (0.6305,0.6004) (0.6508,0.6148) (0.6712,0.6281) (0.6915,0.6404) (0.7119,0.6519) (0.7322,0.6626) (0.7525,0.6726) (0.7729,0.6819) (0.7932,0.6907) (0.8136,0.6989) (0.8339,0.7067) (0.8542,0.7140) (0.8746,0.7210) (0.8949,0.7275) (0.9153,0.7337) (0.9356,0.7396) (0.9559,0.7453) (0.9763,0.7506) (0.9966,0.7557) (1.0169,0.7605) (1.0373,0.7652) (1.0576,0.7696) (1.0780,0.7739) (1.0983,0.7779) (1.1186,0.7818) (1.1390,0.7856) (1.1593,0.7892) (1.1797,0.7926) (1.2000,0.7959)};
\draw[cWalk,densely dashed,thick] plot coordinates {(0.0000,0.0000) (0.0203,0.0035) (0.0407,0.0078) (0.0610,0.0131) (0.0814,0.0196) (0.1017,0.0278) (0.1220,0.0382) (0.1424,0.0516) (0.1627,0.0690) (0.1831,0.0916) (0.2034,0.1210) (0.2237,0.1588) (0.2441,0.2066) (0.2644,0.2655) (0.2847,0.3350) (0.3051,0.4128) (0.3254,0.4939) (0.3458,0.5713) (0.3661,0.6375) (0.3864,0.6865) (0.4068,0.7152) (0.4271,0.7242) (0.4475,0.7167) (0.4678,0.6971) (0.4881,0.6698) (0.5085,0.6386) (0.5288,0.6061) (0.5492,0.5742) (0.5695,0.5438) (0.5898,0.5153) (0.6102,0.4891) (0.6305,0.4650) (0.6508,0.4430) (0.6712,0.4228) (0.6915,0.4043) (0.7119,0.3873) (0.7322,0.3717) (0.7525,0.3572) (0.7729,0.3439) (0.7932,0.3315) (0.8136,0.3200) (0.8339,0.3092) (0.8542,0.2992) (0.8746,0.2897) (0.8949,0.2809) (0.9153,0.2726) (0.9356,0.2647) (0.9559,0.2573) (0.9763,0.2503) (0.9966,0.2436) (1.0169,0.2373) (1.0373,0.2313) (1.0576,0.2256) (1.0780,0.2201) (1.0983,0.2149) (1.1186,0.2100) (1.1390,0.2052) (1.1593,0.2007) (1.1797,0.1963) (1.2000,0.1921)};
\fill[cWalk] (0.4272,0.7242) circle (1.3pt);
\node[anchor=north west,align=left,inner sep=1pt] at (0.43,1.07) {實線 {\color{cBar}$\rho=\langle N\rangle/V$}，虛線 {\color{cWalk}$\mathrm{Var}\,N/V$}\\由淺到深：$L=3,4,5$};
\end{tikzpicture}}

\newcommand{\TenTabRho}{%
\begin{tabular}{@{}l*{6}{c}@{}}
\toprule
$\rho$，$z=$ & 0.2 & $z_c$ & 0.6 & 1.0 & 峰 $z^*$ & 峰高\\
\midrule
$3\times3$ & 0.1391 & 0.3374 & 0.5284 & 0.6823 & 0.4443 & 0.4289\\
$4\times4$ & 0.0900 & 0.3056 & 0.5585 & 0.7338 & 0.4380 & 0.5846\\
$5\times5$ & 0.0603 & 0.2769 & 0.5764 & 0.7565 & 0.4272 & 0.7242\\
\bottomrule
\end{tabular}}

\newcommand{\TenTabExp}{%
\begin{tabular}{@{}lcccc@{}}
\toprule
指數 & Table 10.1 & \shortstack{平均場\\$d=4$} & \shortstack{二維\\$\nu=\frac34$} & \shortstack{三維（引用）\\$\nu\approx0.588$}\\
\midrule
$\beta'=d\nu-1$ & $1-\alpha$ & $1$ & $\frac{1}{2}$ & 0.764\\
$\gamma'=2-d\nu$ & $\alpha$ & $0$ & $\frac{1}{2}$ & 0.236\\
$\alpha'=2-d\nu$ & $\alpha$ & $0$ & $\frac{1}{2}$ & 0.236\\
$\nu'=\nu$ & $\nu$ & $\frac{1}{2}$ & $\frac{3}{4}$ & 0.588\\
\midrule
$\alpha'+2\beta'+\gamma'$ & 2 & 2\,$\checkmark$ & 2\,$\checkmark$ & 2\,$\checkmark$\\
$2-\alpha'$ & & $2$ & $\frac{3}{2}$ & 1.764\\
$d\,\nu'$ & & $2$\,$\checkmark$ & $\frac{3}{2}$\,$\checkmark$ & 1.764\,$\checkmark$\\
書上 $d_c\nu'$ & & $2$ & {\color{cWrong}$3$} & {\color{cWrong}2.352}\\
$d_c=(2-\alpha')/\nu'$ & 4 & 4 & & \\
\bottomrule
\end{tabular}}

\newcommand{\TenTabErrata}{%
\begin{tabular}{@{}lll@{}}
\toprule
位置 & 書上 & 應為\\
\midrule
p.~235，(10.1) & {\color{cWrong}$\sum_NC(N;\vec x,\vec y)\,z^n$} & {\color{cExact}$\sum_NC(N;\vec x,\vec y)\,z^N$}\\
p.~238，(10.14)$\to$(10.15) & {\color{cWrong}對 $S$ 的積分是 $i\pi$；$\Gamma_{MF}=\frac{N_0}4\int_0^\infty f\,d\lambda$} & {\color{cExact}$2\pi i\,\theta(\lambda)$；$\Gamma_{MF}=\frac{N_0}2\int_0^\infty f\,d\lambda$}\\
p.~238，註 3 & {\color{cWrong}$\lambda\to4\lambda$、$4u\to f$、$2r\to1-z/z_c$} & {\color{cExact}$\lambda\to2\lambda$、$u\to f$、$r\to1-z/z_c$}\\
p.~242，(10.32) & {\color{cWrong}$\exp\big[-\frac14\vec h\cdot X^{-1}\cdot\vec h\big]$} & {\color{cExact}$\exp\big[+\frac14\vec h\cdot X^{-1}\cdot\vec h\big]$}\\
p.~242，(10.33) & {\color{cWrong}$\sqrt{\lambda_2/\lambda_1}=\sqrt{\frac{r+q^2+u|\vec M|^2}{r+q^2+3u|\vec M|^2}}$} & {\color{cExact}$\sqrt{\lambda_1/\lambda_2}=\cdots$（右邊對）}\\
p.~245，(10.45) & {\color{cWrong}$e^{-N_0(r\lambda+u\lambda^2)}$} & {\color{cExact}$e^{-N_0(\frac r2\lambda+\frac u4\lambda^2)}$}\\
p.~245，正文 & {\color{cWrong}$D$ \emph{increases} with increasing separation} & {\color{cExact}$-D$（與 $1-D$）隨距離增加；$D$ 遞減}\\
p.~245，(10.46) & {\color{cWrong}$\sum_{\vec q}\frac{\psi_{\vec q}(\vec x)\psi_{\vec q}(\vec x\,')}{q^2(2|r|+q^2)}$} & {\color{cExact}$\sum_{\vec q\ne0}\frac{2|r|\,\psi_{\vec q}(\vec x)\psi_{\vec q}(\vec x\,')}{q^2(2|r|+q^2)}$}\\
pp.~247--248，(10.53)(10.55) & {\color{cWrong}$e^{-\frac V{a^d}\frac{r^2}{4u}}$；$\ln V-\frac V{a^d}\frac{r^2}{4u}$} & {\color{cExact}$e^{+\frac V{a^d}\frac{r^2}{4u}}$；$\ln V+\frac V{a^d}\frac{r^2}{4u}$}\\
p.~251，(10.75) & {\color{cWrong}$2-\alpha'=d_c\nu'$（hold in general）} & {\color{cExact}$2-\alpha'=d\,\nu'$}\\
p.~253，(10.83) & {\color{cWrong}$\mathrm{erfc}\big(\frac{z_cR^2}{2(N+1)}\big)$} & {\color{cExact}$\mathrm{erfc}\big(\sqrt{\frac{z_cR^2}{4(N+1)}}\big)$}\\
\bottomrule
\end{tabular}}


\section{漫步與 \texorpdfstring{O($n$)}{O(n)} 模型：平均場與自旋波（第 10 章）}

\begin{keybox}[先看這裡]
第 9 章說：$n\to0$ 的 O($n$) 磁鐵，關聯函數就是自迴避漫步的生成函數。這一章把磁鐵的兩個標準近似——\textbf{平均場}（所有自旋都一樣）與\textbf{自旋波}（平均場上的小漣漪）——用在\textbf{關在盒子裡的漫步}，
特別是長到把盒子擠滿的「凝聚相」。路線：週期盒子（環面）裡的高分子 $\to$ 平均場化成一維積分，就是第 7 章的 (7.57) $\to$ 自旋波：每個波數 $n-1$ 個橫向、1 個縱向模態，$n\to0$ 時橫向是 $-1$ 個
$\to$ 頭尾之間的\textbf{相關洞}，遠處是 Coulomb 形 $|\vec x-\vec x\,'|^{2-d}$ $\to$ Redner--Reynolds：漫步的\textbf{密度}是序參量，$\beta'=1$、$\gamma'=\alpha'=0$、上臨界維度 4。

\textbf{記號}：步數 $N$、每步的逸度 $z$、$z_c=1/\mu$（二維 0.3791）；這一章 $e^{-r}=z/z_c$（第 9 章是 $e^{-rN/2}$）；O($n$) 的分量數 $n$；格點數 $N_0$（$=V$，格距 1）；
平均場的權重 $f(\lambda)=e^{-N_0(\frac r2\lambda+\frac u4\lambda^2)}$；自旋波的 $\lambda_1$（橫向）、$\lambda_2$（縱向）；相關洞 $D$；漫步的新指數 $\beta',\gamma',\alpha',\nu'$（磁鐵的不加撇號）。
\begin{center}
\small
\begin{tabular}{@{}lll@{\qquad}lll@{}}
\toprule
節 & 內容 & 動畫 & 節 & 內容 & 動畫\\
\midrule
10.1 & 關在盒子裡、環面 & W1001 & 10.3 後半 & 相關洞 $D$、Coulomb 形 & W1004\\
10.2 & 平均場、先取 $n\to0$ & W1002 & 10.3.1 & Redner--Reynolds、密度 & W1005\\
10.3 & 自旋波、負一個模態 & W1003 & 10.4 & 洞的來源：普通漫步的 $1/R$ & W1004\\
\bottomrule
\end{tabular}
\end{center}
（上表是書上的節號；本講義 10.4 節＝書上 10.3 後半＋10.4，10.5 節＝書上 10.3.1，順序和動畫相同。）數字一律由 \texttt{codes/utils/walk\_cond.py} 重算，
逐式校對見 \texttt{ref/md/ch10\_formulas.md}，圖表由 \texttt{ch10/make\_figs.py} 產生並先做數字檢查。連續播放：\texttt{Ch10Demo}（W1001--W1005）。
\end{keybox}

% =====================================================================
\subsection{關在盒子裡的高分子}
\leavevmode\bk{pp.~233--236}\ \anim{W1001}\par\nopagebreak

\textbf{問}：一條高分子關在小盒子裡、和一池單體平衡——盒子的牆會怎樣改變它能走的路？

\textbf{和單體池平衡}（Fig.~10.1）：單體可以接上鏈的兩端，也可以從端點蒸發回池裡，所以鏈長 $N$ 不固定；每接上一個單體得一個因子 $z$（單體的\textbf{逸度}）。從 $\vec x$ 到 $\vec y$ 的鏈的總權重是巨配分函數
\[G(z;\vec x,\vec y)=\sum_NC(N;\vec x,\vec y)\,z^N\tag{10.1}\]
（$C$＝自迴避條數），正是第 9 章的生成函數：$z$ 小時短鏈為主，$z$ 越大鏈越長。
\begin{errbox}[p.~235，(10.1)]
書上印成 $z^n$。$n$ 是 O($n$) 的分量數：$\langle N\rangle=z\,\partial\ln G/\partial z$ 只剩 $n$，$n\to0$ 時 $G\to\sum_NC(N)$ 與 $z$ 無關、$\langle N\rangle=0$（$4\times4$ 環面上永遠是 90677）。應為 $z^N$。
\end{errbox}

\textbf{週期邊界}：真的牆讓靠牆處和中間不一樣（10.3 節的 $\tanh$ 形），所以書上讓走出右邊的從左邊回來、走出上邊的從下邊回來：$L\times L$ 格子黏成\textbf{環面}，格點數 $V=L^2$。
\textbf{算到底：$4\times4$ 環面}（圖 \ref{ten-fig-torus} 上；E、N、W、S＝東、北、西、南，從 $(1,1)$ 出發）：\texttt{WNE}（$N=3$）不越界，和在無限平面上一樣；
\texttt{WSSEENNE}（$N=8$）從下邊出去、由上邊進來，再從上邊出去、由下邊回來；\texttt{WSSEEENNNEEESSW}（$N=15$）越界 3 次，16 格全部走過（Hamilton 路徑）——一盤義大利麵。

\textbf{環面上能走的路比較少}：把環面上的漫步「攤開」到平面（圖 \ref{ten-fig-torus} 上右），得到各點\emph{模 $L$} 互不相同的平面漫步；所以環面的 $c_N$＝平面的 $c_N$ 減去「攤開後有兩點相差 $(kL,lL)\ne(0,0)$」的走法。
相差 $L$ 的倍數至少要走 $L$ 步：$N<L$ 時兩者一樣；$N=L$ 時只扣掉一路直走、繞一圈回到起點的 4 條：$c_4=96=100-4$（$4\times4$）、$c_3=32=36-4$（$3\times3$）、$c_5=280=284-4$（$5\times5$）。
$N$ 再大扣得更多（圖 \ref{ten-fig-torus} 下）：$c_N$ 先增後減，走滿的只剩 5712 條，全部 90677 條（$4\times4$ 環面就是四維超立方體，$16\times5712=91392$ 正是四位元 Gray 碼的個數）。

\textbf{三種情況}：$\langle N\rangle=\sum_NNc_Nz^N\big/\sum_Nc_Nz^N$。$z=0.2$：分母 $1+0.8+0.48+0.288+\cdots=2.8782$，分子 $0.8+0.96+0.864+\cdots=4.1445$，$\langle N\rangle=1.4400$，漫步比盒子小；
$z=z_c=0.3791$：4.8899，和盒子差不多；$z=1.0$：11.7409（最多 15），擠滿盒子。p.~234 強調\textbf{極限的順序}：先取 $n\to0$，再讓盒子變大（10.2 節用數字看為什麼）。

\begin{figure}[htbp]
\centering
\TenFigTorus\\[6pt]
{\footnotesize\setlength{\tabcolsep}{4pt}\TenTabTorus}
\caption{$4\times4$ 環面上的自迴避漫步。上：三條例子（綠點是起點、空心圈是終點；橙色是盒子裡的步，藍色是越界的半步；邊上的同向小箭頭表示對邊黏在一起）與 $N=8$ 那條攤開到平面
（虛線是盒子和它下面的複本）。下：$c_N$ 的比較，$N\le3$（綠）完全相同，$N=4$ 少了 4 條直線}\label{ten-fig-torus}
\end{figure}

\textbf{揭曉}：盒子小、漫步長，漫步就被擠成一盤義大利麵；本章用 O($n$) 模型的平均場與自旋波描述這個「凝聚相」。

% =====================================================================
\subsection{\texorpdfstring{O($n$)}{O(n)} 模型的平均場}
\leavevmode\bk{pp.~236--239}\ \anim{W1002}\par\nopagebreak

\textbf{問}：如果每個自旋都一樣（平均場），$n\to0$ 的 O($n$) 模型會給出什麼？

\textbf{模型}（(10.2)--(10.7)）：$\mathcal H$ 就是第 9 章的 GLW (9.28)（$\frac r2|\vec S_i|^2+\frac12|\nabla\vec S_i|^2+\frac u4|\vec S_i|^4-\vec h\cdot\vec S_i$ 對格點加總），而
\[G(r;\vec R_i-\vec R_j)=\lim_{n\to0}\frac1n\sum_{\alpha=1}^n\langle S_{i\alpha}S_{j\alpha}\rangle=\sum_NC(N;\vec R_i-\vec R_j)\,e^{-rN}.\tag{10.4}\]
對終點加總得磁化率 $\Gamma(r)=\sum_jG=\sum_N\Gamma_Ne^{-rN}$（(10.5)；$\Gamma_N$＝$N$ 步的自迴避條數，以 $z_c^{-N}$ 為單位）。$r>0$ 是 $z<z_c$，$r<0$ 是 $z>z_c$。
〔注〕$e^{-r}=z/z_c$；第 9 章 (9.76)(9.77) 用 $e^{-rN/2}$，兩章的 $r$ 差 2 倍；書上說 $r=(T-T_c)/T_c$。

\textbf{平均場}（(10.8)--(10.11)）：令所有 $\vec S_i=\vec M$，梯度項消失，$Z_{MF}=\int d^nM\,e^{-N_0\{\frac r2|\vec M|^2+\frac u4|\vec M|^4-\vec h\cdot\vec M\}}$，對 $h$ 微分兩次得 $\Gamma_{MF}=N_0\lim_{n\to0}\langle M_1^2\rangle$。
被積函數只依賴 $\lambda=|\vec M|^2$：權重是 $f(\lambda)$，$n$ 維球殼給 $d^nM\propto\lambda^{n/2-1}d\lambda$，對稱性給 $\langle M_1^2\rangle=\langle\lambda\rangle/n$。
$n\to0$ 時分母被 $\lambda\approx0$ 主宰（$\int_0^1\lambda^{n/2-1}d\lambda=\frac2n\to\infty$，分母 $\approx\frac2nf(0)=\frac2n$），兩個 $\frac1n$ 相消：
\[\Gamma_n=\frac{N_0}n\,\frac{\int_0^\infty\lambda^{n/2}f\,d\lambda}{\int_0^\infty\lambda^{n/2-1}f\,d\lambda}\ \xrightarrow{\ n\to0\ }\ \Gamma_{MF}=\frac{N_0}2\int_0^\infty e^{-N_0(\frac r2\lambda+\frac u4\lambda^2)}d\lambda .\tag{10.15}\]
數值（$N_0=100$、$r=0.5$、$u=1$，圖 \ref{ten-fig-mf} 右上）：$n=3,2,1,\allowbreak0.1,0.01$ 得 1.7183、1.7634、1.8131、1.8624、1.8676，趨於 1.8682。
$u=0$（高斯）時任何 $n$ 都是 $\langle M_1^2\rangle=\frac1{N_0r}$、$\Gamma=\frac1r=2$；(10.15) 也給 $\frac{N_0}2\cdot\frac2{N_0r}=\frac1r$。

\textbf{書上的路}（(10.12)--(10.14)）：$\delta(\lambda-|\vec M|^2)=\frac1{2\pi}\int dS\,e^{iS(\lambda-|\vec M|^2)}$ 把 $\vec M$ 的積分變成高斯積分（$M_1^2$ 的平均是 $\frac1{2iS}$；分母 $\to2\pi f(0)=2\pi$，
〔注〕(10.12) 的 $\frac1{2\pi}$ 在分子分母都省掉了），得 $\Gamma_{MF}=\frac{N_0}{2\pi}\iint d\lambda\,dS\,f(\lambda)\frac{e^{iS\lambda}}{2iS}$。
\begin{errbox}[p.~238，(10.14)$\to$(10.15)]
書上說「對 $S$ 的積分就是 $i\pi$」，得 $\Gamma_{MF}=\frac{N_0}4\int_0^\infty f\,d\lambda$。$i\pi$ 是主值；照註 2 的 $-iS\to-iS-\epsilon$，極點在實軸上方，$\int dS\,\frac{e^{iS\lambda}}{S-i\epsilon}=2\pi i\,e^{-\epsilon\lambda}\theta(\lambda)$
（$\epsilon=0.1$：$\lambda=1$ 得 $5.6853i=2\pi ie^{-0.1}$、$\lambda=-1$ 得 0），所以應為 $\frac{N_0}2\int_0^\infty f\,d\lambda$。$u=0$ 時書上的式子得 $\frac1{2r}=1$（應為 2）；數值極限 1.8682，書上 0.9341。
\end{errbox}

\textbf{算到底：閉式}。配方 $\frac r2\lambda+\frac u4\lambda^2=\frac u4(\lambda+\frac ru)^2-\frac{r^2}{4u}$ 得
$I\equiv\int_0^\infty f\,d\lambda=\sqrt{\frac\pi{N_0u}}\,e^{N_0r^2/4u}\,\mathrm{erfc}\big(\frac r2\sqrt{N_0/u}\big)$。$N_0=100$、$r=0.5$、$u=1$：$I=0.1772\times518.01\times4.070\times10^{-4}=0.03736$，$\Gamma_{MF}=50I=1.8682$。
用 $\mathrm{erfc}(x)\approx\frac{e^{-x^2}}{x\sqrt\pi}$（$x\to+\infty$）與 $\mathrm{erfc}(-\infty)=2$ 得兩種漸近（圖 \ref{ten-fig-mf}）：
\begin{itemize}[nosep]
\item $r>0$：$f$ 從 $\lambda=0$ 就遞減（$\vec M\approx0$），$\Gamma_{MF}\approx\frac1r$，和普通漫步的 $\sum_Ne^{-rN}\approx\frac1r$ 一樣：漫步比盒子小。$N_0=100,\allowbreak1000,\allowbreak10^4$：1.8682、1.9844、1.9984 $\to2$。
\item $r<0$：$f$ 在 $\lambda^*=|r|/u$ 有尖峰，$\Gamma_{MF}\approx\sqrt{\pi N_0/u}\,e^{N_0r^2/4u}$ 隨體積\emph{指數}增長：漫步擠滿盒子。精確 $\div$ 漸近 0.9998、1.0000、1.0000；
$\ln\Gamma_{MF}/N_0=0.0912,\allowbreak\ 0.0665,\allowbreak\ 0.0630\to\frac{r^2}{4u}=0.0625$（p.~239 的「reduction factor」$e^{Nr^2/4u}$ 就是這個鞍點值）。
\end{itemize}
\textbf{為什麼要先取 $n\to0$}：$r=-0.5$ 時 $\Gamma_{MF}=9179.7$，而 $n=10^{-2},\allowbreak10^{-3},\allowbreak10^{-4},\allowbreak10^{-6}$ 的 $\Gamma_n/\Gamma_{MF}=0.3299,\allowbreak\ 0.8311,\allowbreak\ 0.9801,\allowbreak\ 0.9998$。推導要分母的 $\frac2n$ 壓過 $\lambda^*$ 附近的尖峰
（$\propto e^{N_0r^2/4u}=518.01$），即 $n\,e^{N_0r^2/4u}\ll1$：盒子越大 $n$ 要越小，先讓盒子變大就到不了 $n\to0$ 的結果。

\textbf{和第 7 章對照}（註 3）：在 (10.15) 裡換 $N_0\to N$、$\lambda\to2\lambda$、$u\to f$、$r\to1-x$（$x=z/z_c$），得 $\Gamma_{MF}=N\int_0^\infty e^{N(-f\lambda^2-\lambda(1-x))}d\lambda=N\,G_{MF}$，
$G_{MF}$ 就是第 7 章的平均場 (7.57)；多出的 $N$ 是對終點加總。$r=0.5$（$x=\frac12$，第 7 章練習 7.5 的 $x$）：$\Gamma_{MF}=1.8682=100\,G_{MF}(\frac12)$。$r\to1-z/z_c$ 正是 $e^{-r}=z/z_c$ 在 $z_c$ 附近的展開。
\begin{errbox}[p.~238，註 3]
書上的對照 $\lambda\to4\lambda$、$4u\to f$、$2r\to1-z/z_c$ 只對它自己的 $\frac{N_0}4$ 成立（兩個 2 倍抵掉：$\frac{N_0}4I=0.9341$，照書上的換法 $100\,G_{MF}$ 也是 0.9341）。
配上更正的 (10.15)，應為 $\lambda\to2\lambda$、$u\to f$、$r\to1-z/z_c$（也和 (10.4) 的 $e^{-r}=z/z_c$ 一致）。
\end{errbox}

\begin{figure}[htbp]
\centering
\begin{minipage}[c]{0.42\linewidth}\centering\TenFigMF\end{minipage}\hfill
\begin{minipage}[c]{0.56\linewidth}\centering\footnotesize\setlength{\tabcolsep}{3pt}\TenTabN\\[8pt]\TenTabAsym\end{minipage}
\caption{平均場（$N_0=100$、$u=1$）。左：$f(\lambda)$ 除以峰值，$r=0.5$（藍）從 0 遞減、$r=-0.5$（橙）在 $\lambda^*$ 有尖峰。
右上：$\Gamma_n=N_0\langle M_1^2\rangle_n$ 趨於 (10.15)，$r<0$ 時 $n$ 要非常小。右下：$N_0$ 變大時的兩種漸近}\label{ten-fig-mf}
\end{figure}

\textbf{揭曉}：平均場把盒子裡的漫步變成一維積分；$r<0$ 多出的 $e^{N_0r^2/4u}$ 正是第 7 章的因子。

% =====================================================================
\subsection{自旋波：在平均場上加漣漪}
\leavevmode\bk{pp.~239--243}\ \anim{W1003}\par\nopagebreak

\textbf{問}：平均場把所有自旋看成一樣——在它上面加一點漣漪（自旋波），會多出什麼？

\textbf{狀態方程}（(10.16)(10.17)）：連續版的 $\mathcal H$ 把 $\sum_i$ 換成 $\int d^dx$；讓它極小的組態 $\vec M(\vec x)$ 滿足下面左邊的狀態方程（練習 10.1）。\textbf{算到底：牆邊}（圖 \ref{ten-fig-sw} 左；書上沒寫的例子）：
一維、$h=0$、牆上 $M(0)=0$、遠處 $M\to M_0=\sqrt{|r|/u}$（$r<0$）。試 $M=M_0\tanh kx$：$M''=-2k^2M(1-M^2/M_0^2)$，而 (10.17) 要 $M''=rM+uM^3=-|r|M(1-M^2/M_0^2)$，所以 $k=\sqrt{|r|/2}$：
\[-\nabla^2\vec M+r\vec M+u\vec M|\vec M|^2-\vec h=0,\qquad M(x)=\sqrt{|r|/u}\,\tanh\big(x\sqrt{|r|/2}\big).\tag{10.17}\]
$r=-1$、$u=1$：$M(1),M(2),M(4)=0.6089,\allowbreak\ 0.8884,\allowbreak\ 0.9930$；牆的影響只延伸約一個\textbf{癒合長度} $\sqrt{2/|r|}=1.4142$（有限差分＋牛頓法解邊值問題，和 $\tanh$ 只差 $2.8\times10^{-6}$）。
週期邊界沒有牆，$\vec M$ 就是常數——以下都用週期邊界。

\textbf{展開到二次}（(10.18)--(10.21)）：令 $\vec S=\vec M+\vec\sigma$，把 $|\vec M+\vec\sigma|^4=(|\vec M|^2+2\vec M\cdot\vec\sigma+|\vec\sigma|^2)^2$ 乘開。一次項是 (10.17) 的左邊乘 $\vec\sigma$，等於 0（極小點附近沒有一次項）；二次項
\[\mathcal H_2=\int d^dx\Big\{\frac12|\nabla\vec\sigma|^2+\frac r2|\vec\sigma|^2+\frac u2|\vec M|^2|\vec\sigma|^2+u(\vec M\cdot\vec\sigma)^2\Big\},\tag{10.21}\]
最後一項只罰「沿著 $\vec M$」的漣漪——這就是橫向與縱向的差別。週期盒子裡用 Fourier 模態 $\psi_{\vec q}=\prod_ie^{iq_ix_i}/\sqrt{L_i}$、$q_i=2\pi n_i/L_i$（每個方向剛好繞整數圈；(10.22)--(10.24)），
$\nabla$ 變成乘 $i\vec q$，各個 $\vec q$ 分開：$\mathcal H_2=\sum_{\vec q}\vec\sigma_{\vec q}^*\cdot X_{\vec q}\cdot\vec\sigma_{\vec q}$，
\[X_{\vec q}=\lambda_1I+u|\vec M\rangle\langle\vec M|,\qquad\lambda_1=\frac r2+\frac12q^2+\frac12u|\vec M|^2,\qquad\lambda_2=\lambda_1+u|\vec M|^2.\tag{10.27--10.29}\]
垂直於 $\vec M$ 的向量被 $|\vec M\rangle\langle\vec M|$ 打成 0，特徵值 $\lambda_1$（$n-1$ 重，\textbf{橫向}）；沿著 $\vec M$ 多 $u|\vec M|^2$，特徵值 $\lambda_2$（1 重，\textbf{縱向}；練習 10.3）。

\textbf{算到底：$n=3$}。$r=-1$、$u=1$、$q^2=1$、$\vec M=(\frac23,\frac23,\frac13)$（$|\vec M|^2=1=|r|/u$，平均場的鞍點）：$\lambda_1=-\frac12+\frac12+\frac12=\frac12$、$\lambda_2=\frac32$；
(10.30) 給 $X^{-1}=\frac I{\lambda_1}-\frac{u|\vec M\rangle\langle\vec M|}{\lambda_1\lambda_2}=2I-\frac43|\vec M\rangle\langle\vec M|$：
\[X=\frac12I+|\vec M\rangle\langle\vec M|=\frac1{18}\begin{pmatrix}17&8&4\\8&17&4\\4&4&11\end{pmatrix},\qquad X^{-1}=\frac1{27}\begin{pmatrix}38&-16&-8\\-16&38&-8\\-8&-8&50\end{pmatrix}.\tag{10.30}\]
核對：$X\vec M=\frac32\vec M$；$(1,-1,0)$、$(1,1,-4)$ 垂直於 $\vec M$，都乘 $\frac12$；跡 $\frac{45}{18}=\frac12+\frac12+\frac32$；$XX^{-1}=X^{-1}X=I$（分數精確，(10.31)）；隨機的 $\vec M$、$n=1$ 到 6 誤差 $<10^{-11}$。
$q=0$ 時 $\lambda_1=0$：整體轉動 $\vec M$ 不花能量。

\textbf{高斯積分}（(10.32)）：每個 $\vec q$ 是 $n$ 維高斯積分 $\int d^n\sigma\,e^{-\vec\sigma\cdot X\cdot\vec\sigma+\vec h\cdot\vec\sigma}=\frac{\pi^{n/2}}{\sqrt{\det X}}\,e^{+\frac14\vec h\cdot X^{-1}\cdot\vec h}$，$\det X=\lambda_1^{n-1}\lambda_2$。
例：$n=2$、$\vec M=(\frac35,\frac45)$、$\vec h=(\frac12,0)$：$\det X=\frac34$，$(X^{-1})_{11}=2-\frac43\cdot\frac9{25}=\frac{38}{25}$，$\vec h\cdot X^{-1}\cdot\vec h=0.38$；積分 $=\frac{\pi}{\sqrt{3/4}}\,e^{0.38/4}=3.6276\times e^{0.095}=3.9891$，和數值積分相同。
\begin{errbox}[p.~242，(10.32)(10.33)]
(1) (10.32) 的指數印成 $\exp[-\frac14\vec h\cdots X^{-1}\cdots\vec h]$，應為 $+\frac14$（配方：$-\vec\sigma\cdot X\cdot\vec\sigma+\vec h\cdot\vec\sigma$ 的極大值是 $+\frac14\vec h\cdot X^{-1}\cdot\vec h$）；上例照書上得 3.2988。
(2) (10.33) 左邊寫 $\sqrt{\lambda_2/\lambda_1}$，應為 $\sqrt{\lambda_1/\lambda_2}$（右邊的 $\sqrt{\frac{r+q^2+u|\vec M|^2}{r+q^2+3u|\vec M|^2}}$ 對）；$n=3$ 的例子照書上得 $\sqrt3=1.7321$，應為 0.5774。
\end{errbox}
\textbf{負一個模態}（(10.33)，圖 \ref{ten-fig-sw} 右）：$(\det X)^{-1/2}=\lambda_1^{-(n-1)/2}\lambda_2^{-1/2}$。$n=3$ 兩個橫向、一個縱向；$n=1$（Ising）只有縱向；$n\to0$ 時橫向的個數 $n-1\to-1$，
$\frac1{\sqrt{\lambda_2}\sqrt{\lambda_1^{n-1}}}\to\sqrt{\frac{\lambda_1}{\lambda_2}}=\sqrt{\frac{r+q^2+u|\vec M|^2}{r+q^2+3u|\vec M|^2}}$——「負一個模態」把 $\sqrt{\lambda_1}$ 搬到分子：數學上允許，物理上奇妙。上例 $\to\frac1{\sqrt3}=0.5774$。

\begin{figure}[htbp]
\centering
\begin{minipage}[c]{0.5\linewidth}\centering\TenFigWall\end{minipage}\hfill
\begin{minipage}[c]{0.48\linewidth}\centering\footnotesize\setlength{\tabcolsep}{3pt}\TenTabModes\end{minipage}
\caption{自旋波。左：牆邊的狀態方程解（$r=-1$、$u=1$），切線（點線）在癒合長度碰到 $M_0$。
右：$\vec M=(\frac23,\frac23,\frac13)$、$q^2=1$ 的模態；$n\to0$ 時橫向模態的個數是 $-1$，$(\det X)^{-1/2}\to\sqrt{\lambda_1/\lambda_2}=1/\sqrt3$}\label{ten-fig-sw}
\end{figure}

\textbf{揭曉}：每個 $\vec q$ 有 $n-1$ 個橫向、1 個縱向模態；$n\to0$ 時橫向變成 $-1$ 個，留下 $\sqrt{\lambda_1/\lambda_2}$。

% =====================================================================
\subsection{相關洞：頭尾互相排斥}
\leavevmode\bk{pp.~243--246、251--253}\ \anim{W1004}\par\nopagebreak

\textbf{問}：擠滿盒子的漫步，頭和尾喜不喜歡靠在一起？

\textbf{生成函數}（(10.37)）：$G(z;\vec x,\vec x\,')=\lim_{n\to0}\frac1n\frac{\partial^2\ln Z}{\partial h(\vec x)\partial h(\vec x\,')}$——外場隨位置變，兩次微分取出「從 $\vec x$ 到 $\vec x\,'$」的漫步。書上 (10.32)--(10.43) 的長式只看結構，分四步：
\begin{enumerate}[label=\textbf{\arabic*.},itemsep=0em,leftmargin=1.6em]
\item \textbf{對漣漪積分}（(10.32)(10.33)）：每個 $\vec q$ 一個高斯積分；$n\to0$ 時行列式剩 $\prod_{\vec q}\sqrt{\lambda_1/\lambda_2}$，外場留在 $e^{+\frac14\vec h\cdot X^{-1}\cdot\vec h}$ 裡。
\item \textbf{把 $X^{-1}$ 拆兩塊}（(10.34)--(10.36)）：與 $\vec M$ 的方向無關的 $\vec h\cdot\vec h\,'$ 項（係數 $G_0$），與沿著 $\vec M$ 的 $(\vec h\cdot\vec M)(\vec h\,'\cdot\vec M)$ 項——係數是橫向與縱向「傳播子」之差
（$\vec q=0$ 的零模是 $\vec M$ 本身的轉動，要去掉）：
\[D(\vec x,\vec x\,')=\sum_{\vec q\ne0}\psi_{\vec q}(\vec x)\psi_{\vec q}(\vec x\,')\Big(\frac1{r+q^2+u|\vec M|^2}-\frac1{r+q^2+3u|\vec M|^2}\Big).\tag{10.36}\]
\item \textbf{對 $\vec M$ 積分}（(10.38)--(10.42)）：用 10.2 節的 $\lambda$ 技巧，$n\to0$ 時 $\int f\,d^nM\to1$、$\int(\vec h\cdot\vec M)(\vec h\,'\cdot\vec M)f\,d^nM\to\frac{\vec h\cdot\vec h\,'}2\int_0^\infty f\,d\lambda$（就是 (10.15) 的 $\frac{N_0}2$）。
\item \textbf{取係數}（(10.43)）：$\vec h\cdot\vec h\,'$ 的係數是 $\frac14\{G_0-(\int_0^\infty f\,d\lambda)\,D\}$。〔注〕本系列只核對到這一層結構，沒有逐行驗算。
\end{enumerate}

\textbf{兩種情況}：$r\gg0$ 時 $\vec M=0$，$G_0\to\int\frac{d^dq}{(2\pi)^d}\frac{e^{i\vec q\cdot(\vec x-\vec x\,')}}{r+q^2}$（(10.44)）：看不到盒子的普通漫步（$32^3$ 週期格子、$r=1$，和無限大格子相差 $1.2\times10^{-11}$）。
$r\ll0$ 時平均場的峰主宰；在鞍點 $u|\vec M|^2=|r|$，(10.36) 的括號是 $\frac1{q^2}-\frac1{q^2+2|r|}$，於是 $G$ 是平均場的條數 $\times(1-D)$：
\[G\propto\int_0^\infty e^{-N_0(\frac r2\lambda+\frac u4\lambda^2)}d\lambda\,\big(1-D\big),\qquad D(\vec x,\vec x\,')=\sum_{\vec q\ne0}\frac{2|r|\,\psi_{\vec q}(\vec x)\psi_{\vec q}(\vec x\,')}{q^2(2|r|+q^2)} .\tag{10.45, 10.46}\]
\textbf{算到底：$64^3$ 週期格子}（$r=-1$、$u=1$，格子上 $q^2\to\sum_i4\sin^2\frac{q_i}2$；圖 \ref{ten-fig-hole} 左）：沿軸相距 $\Delta=0,1,2,4,\allowbreak8,16,32$，
$D=0.1083,\allowbreak\ 0.0613,\allowbreak\ 0.0359,\allowbreak\ 0.0167,\allowbreak\ 0.0065,\allowbreak\ 0.0016,\allowbreak\ -0.0001$，$1-D=0.8917\to1.0001$：頭尾重合時條數少 11\%、相鄰時少 6\%，越遠越不受影響——這個凹陷就是\textbf{相關洞}。
$-D\approx c_0-c/\Delta$（$\Delta=4$ 到 16 擬合）：$c=0.0801$，和 $\frac1{4\pi}=0.0796$ 只差 1.0064 倍。
\begin{errbox}[p.~245，(10.45)(10.46) 與正文]
(1) (10.45) 的指數印成 $-N_0(r\lambda+u\lambda^2)$，應為 (10.15) 的 $-N_0(\frac r2\lambda+\frac u4\lambda^2)$：照書上極大在 $|r|/2u=0.5$，不是正文說的 $|r|/u=1$（$r=-1$）。
(2) 正文說 $D$「隨距離增加」：$D$ 的權重都是正的，$D(0)$ 最大、隨距離遞減（$\Delta=0\to8$：0.1083 $\to$ 0.0065）；增加的是 $-D$（書上 Fig.~10.2 畫的就是 $-D$）與 $1-D$。
(3) (10.46) 的分子少了 $2|r|$，也要去掉 $\vec q=0$：照書上 $D(0)=0.0541$，只有一半（(10.47) 寫「正比於」，不受影響）。
\end{errbox}

\textbf{Coulomb 形}（(10.47)(10.48)、練習 10.5）：三維、無限大的連續版積得出來：$\frac1{q^2}$ 與 $\frac1{q^2+2|r|}$ 的 Fourier 轉換是 Coulomb 與 Yukawa，
$D(\rho)=\frac1{4\pi\rho}(1-e^{-\sqrt{2|r|}\rho})\to\frac1{4\pi\rho}$（$\rho\gg1/\sqrt{2|r|}$），與 $r$ 無關；一般維度 $D\propto|\vec x-\vec x\,'|^{2-d}$（(10.48)）。
$r=-1$、$\rho=1,2,4,8$：0.0602、0.0374、0.0198、0.0099；無限大格子 0.0648、0.0394、0.0202、0.0100（圖 \ref{ten-fig-hole} 右上）；週期盒子只讓 $D$ 整體下移一個常數
$\frac{2.837297}{4\pi L}=0.0035$（簡單立方的 Ewald 常數，就是上面的 $c_0$）。〔注〕p.~246 的條件寫成 $\gg1/|r|$；長度尺度是 $1/\sqrt{2|r|}$。

\textbf{洞從哪裡來}（書上 10.4 節，(10.76)--(10.84)）：先看\emph{普通}漫步。從原點走 $N$ 步，落在離起點 $\vec R$ 處的步數期望值是 $\sum_{n\le N}P_n(\vec R)$。
Fourier 空間裡是公比 $\frac{z_c}{z_c+k^2}$ 的等比級數，加到第 $N$ 步得 $\frac1{z_ck^2}[1-e^{-(N+1)k^2/z_c}]$（(10.80)(10.81)），轉回實空間：
\[\sum_{n\le N}P_n(R)\ \propto\ \frac1{4\pi z_cR}\,\mathrm{erfc}\Big(\sqrt{\frac{z_cR^2}{4(N+1)}}\Big)\tag{10.83}\]
——$R$ 比 $2\sqrt{N/z_c}$ 小時 $\propto\frac1R$，再遠就被高斯截斷。\textbf{核對}（簡單立方格子、每步 $\frac16$，$z_c\leftrightarrow6$、前因子 $\frac3{2\pi R}$；$N=100$；圖 \ref{ten-fig-hole} 右下）：
精確計數和 erfc 形只差 3--11\%，轉折在 $R\approx8.2$；$N\to\infty$ 時 $G(R)\cdot\frac{2\pi R}3=1.0815\to1.0009$，就是 $\frac1R$。
所以長漫步在自己任一點的周圍，步數的密度 $\propto R^{-(d-2)}$；自迴避讓頭（或尾）在這些地方受罰，修正 $\sim-u/R$，和 (10.48) 的洞是同一個形狀。
\begin{errbox}[p.~253，(10.83)]
erfc 的引數印成 $\frac{z_cR^2}{2(N+1)}$：$\int\frac{d^3k}{(2\pi)^3}\frac{e^{i\vec k\cdot\vec R}}{k^2}e^{-ak^2}=\frac{\mathrm{erf}(R/2\sqrt a)}{4\pi R}$、$a=\frac{N+1}{z_c}$，應為 $\sqrt{\frac{z_cR^2}{4(N+1)}}$。
照書上 $R=8$ 得 $4.3\times10^{-4}$（精確 0.0103）、$R=16$ 得 $1.7\times10^{-28}$：尾巴是 $e^{-(z_cR^2/2N)^2}$，不是高斯。
\end{errbox}
〔注〕(10.79) 的 $d^kk_1$ 應為 $d^dk_1$；(10.80) 書上省了 $z_c^{-N}$；p.~253 的轉折 $R\le\sqrt{2N/z_c}$ 更正後是 $R\approx2\sqrt{N/z_c}$；前因子依正規化而定，只看形狀。

\begin{figure}[htbp]
\centering
\begin{minipage}[c]{0.46\linewidth}\centering\TenFigHole\end{minipage}\hfill
\begin{minipage}[c]{0.52\linewidth}\centering\footnotesize\setlength{\tabcolsep}{2.5pt}\TenTabHole\\[8pt]\TenTabVisits\end{minipage}
\caption{相關洞。左：$64^3$ 週期格子沿軸的 $-D(\Delta)$（$r=-1$；書上 Fig.~10.2 的樣子）與 Coulomb 擬合。右上：各種 $D$ 與 $\frac1{4\pi\rho}$。
右下：簡單立方格子普通漫步「離起點 $R$ 的步數」（$N=100$）與更正的 (10.83)，紅字是書上的引數，最後一列是 $N=\infty$}\label{ten-fig-hole}
\end{figure}

\textbf{揭曉}：普通漫步在自己周圍的步數密度 $\propto1/R$；自迴避的懲罰 $-u/R$ 讓頭尾互相排斥——一個 Coulomb 形的洞。

% =====================================================================
\subsection{Redner--Reynolds：漫步的密度是序參量}
\leavevmode\bk{pp.~246--251}\ \anim{W1005}\par\nopagebreak

\textbf{問}：把漫步的步數當成粒子數——它的「密度」能不能像磁化一樣當序參量？

\textbf{巨正則系綜}（(10.49)--(10.52)、(10.59)）：$\Gamma(r)=\sum_N\Gamma_Ne^{-rN}$ 就像氣體的巨配分函數（$N$＝粒子數、$-r$＝化學勢）：
\[\Omega=\frac1V\ln\Gamma,\qquad\langle N\rangle=-\frac{\partial\ln\Gamma}{\partial r},\qquad\rho=\frac{\langle N\rangle}V=-\frac{\partial\Omega}{\partial r},\qquad\frac{\partial^2\Omega}{\partial r^2}=\frac{\mathrm{Var}\,N}V .\tag{10.49--10.59}\]
$e^{-r}=z/z_c$，所以 $-\partial/\partial r=z\,\partial/\partial z$：和 10.1 節的 $\langle N\rangle$ 相同。
\textbf{算到底：環面}（圖 \ref{ten-fig-rr} 左）：$\rho(z)$ 從 0 爬到 $\frac{V-1}V$（$z\to\infty$ 只剩走滿的漫步：$\frac89,\frac{15}{16},\frac{24}{25}$）。
漲落 $\mathrm{Var}\,N/V$ 的峰在 $z^*=0.4443,\allowbreak\ 0.4380,\allowbreak\ 0.4272$，高 0.4289、0.5846、0.7242（$L=3,4,5$）——格子越大峰越高、越往 $z_c$ 靠：有限大小的相變。
無限大的盒子裡，$\rho$ 在 $z_c$ 以下是 0、以上不是 0：密度就是序參量。

\textbf{平均場}（(10.53)--(10.64)）：(10.15) 再對起點加總（多一個 $V$）。$r>0$：$\Gamma_{MF}\propto V/r$，$\Omega=\frac1V\{\ln V-\ln r+\text{常數}\}$，$\rho_{MF}=\frac1{Vr}\to0$；
$r<0$：$\Gamma_{MF}\propto V^{3/2}e^{+Vr^2/4u}$，$\Omega=\frac1V\{\frac32\ln V+\frac{Vr^2}{4u}+\text{常數}\}$，$\rho_{MF}=\frac{|r|}{2u}$（(10.55)--(10.58)）。
$V=10^6$、$u=1$：$r=-0.5,\allowbreak-0.1,\allowbreak-0.01$ 得 $\rho=0.2500,\allowbreak\ 0.0500,\allowbreak\ 0.0050$；$r=0.1,\allowbreak\ 0.5$ 得 $1.0\times10^{-5}$、$2.0\times10^{-6}$。$\ln\rho$ 對 $\ln|r|$ 的斜率 1.0000，所以 $\beta'=1$（磁鐵的 $\beta=\frac12$：$\rho=\frac12|\vec M|^2$，$\beta'=2\beta$）。
$\chi'=\partial^2\Omega/\partial r^2$：$r<0$ 時 $\frac1{2u}=0.5000$；$r=0.1,\allowbreak\ 0.5$ 時 $1.0\times10^{-4}$、$4.0\times10^{-6}$（$=\frac1{Vr^2}\to0$）——只跳一下、不發散：$\gamma'=\alpha'=0$；另有 $\nu'=\frac12$（(10.62)--(10.64)）。
\begin{errbox}[pp.~247--248，(10.53)(10.55)]
$r<0$ 的指數印成 $e^{-\frac V{a^d}\frac{r^2}{4u}}$，$\Omega$ 裡是 $\ln V-\frac V{a^d}\frac{r^2}{4u}$；(10.15) 在 $\lambda^*=|r|/u$ 的值是 $e^{+N_0r^2/4u}$，應為 $+$。照書上的負號 $\rho=-\partial\Omega/\partial r=\frac r{2u}<0$（$r=-0.5$ 得 $-0.2500$）。
\end{errbox}
〔注〕(10.53) 是 $N_0\times$(10.15)，沿用書上的 $\frac{N_0}4$，更正後乘 2（常數，不影響 $\rho$）；(10.55) 的 $\ln V$ 應為 $\frac32\ln V$；(10.58) 的「$\rho_{MF}\to|r|$」是 $\propto|r|$；p.~249 的 $\gamma'_{MF}=0$ 見 (10.63)。

\textbf{標度}（(10.65)--(10.73)）：仿照第 9 章，$\Omega$ 的奇異部分只透過 $rL^{1/\nu}$ 依賴 $r$：$\Omega\sim L^{-d}H(rL^{1/\nu})$（(10.70)），所以 $\rho=-\partial\Omega/\partial r=-L^{1/\nu-d}H'(rL^{1/\nu})$（(10.71)）。
$L\to\infty$ 時 $\rho$ 不能依賴 $L$，只能 $H'(y)\propto|y|^{\nu d-1}$：$\rho\propto|r|^{\nu d-1}$、$\beta'=d\nu-1$（(10.72)(10.73)）；同理 $\alpha'=\gamma'=2-d\nu$、$\nu'=\nu$。平均場（$d=4$、$\nu=\frac12$）回到 $\beta'=1$、$\gamma'=\alpha'=0$。
標度律 $\alpha'+2\beta'+\gamma'=(2-d\nu)+2(d\nu-1)+(2-d\nu)=2$；平均場的 $2-\alpha'=d_c\nu'$ 給 $2=d_c\cdot\frac12$，上臨界維度 $d_c=4$（(10.65)(10.66)）。
〔注〕(10.71) 書上 $\rho=\partial\Omega/\partial r$ 少了負號；(10.72) 的 $r^{\nu d-1}$ 應為 $|r|^{\nu d-1}$。
\textbf{Table 10.1}（p.~251，圖 \ref{ten-fig-rr} 右）用第 9 章的 hyperscaling $d\nu=2-\alpha$ 改寫：$\beta'=1-\alpha$、$\gamma'=\alpha'=\alpha$、$\nu'=\nu$、$d_c=4$——漫步的新指數全由磁鐵的 $\alpha$、$\nu$ 決定
（二維 $\nu=\frac34$：$\beta'=\gamma'=\alpha'=\frac12$；三維 $\nu\approx0.588$ 只引用：$\beta'=0.764$、$\gamma'=\alpha'=0.236$）。
\begin{errbox}[p.~251，(10.75)]
書上寫 $2-\alpha'=d_c\nu'$，並說兩條標度律「一般都成立」。一般維度要用 $d$：$2-\alpha'=d\,\nu'$（就是 hyperscaling）。$d=2$：$2-\alpha'=\frac32=d\nu'$，但 $d_c\nu'=3$；$d=3$：1.764 對 2.352。
\end{errbox}

\begin{figure}[htbp]
\centering
\begin{minipage}[c]{0.46\linewidth}\centering\TenFigRho\\[6pt]\footnotesize\setlength{\tabcolsep}{2.5pt}\TenTabRho\end{minipage}\hfill
\begin{minipage}[c]{0.52\linewidth}\centering\footnotesize\setlength{\tabcolsep}{3pt}\TenTabExp\end{minipage}
\caption{Redner--Reynolds。左：環面的精確 $\rho(z)$ 與漲落 $\mathrm{Var}\,N/V$（圓點是峰），格子越大峰越高、越靠近 $z_c$。
右：漫步的新指數與標度律的驗算；$d\nu'$ 是更正的 (10.75)，紅字是照書上的 $d_c\nu'$}\label{ten-fig-rr}
\end{figure}

\textbf{揭曉}：漫步的密度在 $z_c$ 兩邊從 0 變成有限值——擠滿盒子的漫步就是 O($n$) 磁鐵的有序相。

% =====================================================================
\subsection{練習 10.1--10.5 解答}
{\small
\begin{enumerate}[label=\textbf{10.\arabic*},itemsep=0.05em]
\item 由 $\delta\mathcal H=0$ 推出 (10.17)。\quad 答：$\vec S\to\vec S+\delta\vec S$：$\delta\mathcal H=\int d^dx\{(r\vec S+u\vec S|\vec S|^2-\vec h)\cdot\delta\vec S+\nabla\vec S\cdot\nabla\delta\vec S\}$；
最後一項分部積分成 $-\nabla^2\vec S\cdot\delta\vec S$（週期邊界，或牆上 $\delta\vec S=0$，邊界項為 0）。對任何 $\delta\vec S$ 都要 $\delta\mathcal H=0$ $\Rightarrow-\nabla^2\vec M+r\vec M+u\vec M|\vec M|^2-\vec h=0$。
離散化 $\mathcal H$ 的數值梯度＝格距 $\times$(10.17) 的左邊（相對誤差 $<10^{-7}$）；牆邊的解見 10.3 節。
\item 展開的三階項。\quad 答：$(|\vec M|^2+2\vec M\cdot\vec\sigma+|\vec\sigma|^2)^2$ 的三階項是 $2\cdot(2\vec M\cdot\vec\sigma)|\vec\sigma|^2$，乘 $\frac u4$：$\mathcal H_3=u\int d^dx\,(\vec M\cdot\vec\sigma)|\vec\sigma|^2$；四階 $\frac u4\int|\vec\sigma|^4$；
梯度項與 $r$ 項只到二階。$|\vec M+\epsilon\vec\sigma|^4$ 對 $\epsilon$ 的係數（數值內插）＝公式，誤差 $<10^{-11}$。
\item $X$ 的特徵值。\quad 答：$X=\lambda_1I+uP$、$P=|\vec M\rangle\langle\vec M|$。$P\vec M=|\vec M|^2\vec M$ $\Rightarrow$ 沿 $\vec M$ 的特徵值 $\lambda_1+u|\vec M|^2=\lambda_2$（1 重）；$\vec v\perp\vec M$ $\Rightarrow$ $P\vec v=0$，特徵值 $\lambda_1$（$n-1$ 維）；
所以 $\det X=\lambda_1^{n-1}\lambda_2$（例子 $\frac12,\frac12,\frac32$）。
\item 驗證 (10.30)(10.31)。\quad 答：用 $P^2=|\vec M|^2P$：$(\lambda_1I+uP)\big(\frac I{\lambda_1}-\frac{uP}{\lambda_1\lambda_2}\big)=I+uP\,\frac{\lambda_2-\lambda_1-u|\vec M|^2}{\lambda_1\lambda_2}=I$；$P$ 與 $I$ 可交換，左逆就是右逆（例子見 10.3 節，分數精確）。
\item $D$ 的長距離行為 (10.48)。\quad 答：令 $\vec q=\vec p/|\vec\Delta|$（$\vec\Delta=\vec x-\vec x\,'$）：
\[\int d^dq\frac{e^{i\vec q\cdot\vec\Delta}}{q^2(2|r|+q^2)}=|\Delta|^{4-d}\int d^dp\frac{e^{i\vec p\cdot\hat\Delta}}{p^2(2|r||\Delta|^2+p^2)}\approx\frac{|\Delta|^{2-d}}{2|r|}\int d^dp\frac{e^{i\vec p\cdot\hat\Delta}}{p^2}\qquad(|\Delta|^2|r|\gg1),\]
即 $D\propto|\Delta|^{2-d}$；乘上 (10.46) 漏掉的 $2|r|$ 後與 $r$ 無關。三維精確 $D=\frac1{4\pi\rho}(1-e^{-\sqrt{2|r|}\rho})$：0.0602、0.0374、0.0198、0.0099（$r=-1$）；格子與週期盒子見圖 \ref{ten-fig-hole}。
\end{enumerate}}

% =====================================================================
\subsection{第 10 章的回顧與筆誤一覽}
\begin{center}
\footnotesize\setlength{\tabcolsep}{4pt}
\begin{tabular}{@{}lll@{}}
\toprule
主題（書上的式子） & 結果 & 動畫\\
\midrule
盒子裡的漫步 (10.1) & $4\times4$ 環面的 $c_N$ 從 $N=4$ 起變少（96 對 100），走滿的 5712 條 & W1001\\
平均場 (10.2)--(10.15) & $\Gamma_{MF}=\frac{N_0}2\int_0^\infty f\,d\lambda$；$r<0$ 多出 $e^{N_0r^2/4u}$；先取 $n\to0$ & W1002\\
自旋波 (10.16)--(10.36) & $\lambda_1$（$n-1$ 重）、$\lambda_2$（1 重）；$n\to0$ 剩 $\sqrt{\lambda_1/\lambda_2}$ & W1003\\
相關洞 (10.37)--(10.48)、(10.76)--(10.84) & $1-D$：0.8917 $\to$ 1；$D\approx\frac1{4\pi\Delta}$；步數密度 $\propto1/R$ & W1004\\
Redner--Reynolds (10.49)--(10.75) & $\rho_{MF}=\frac{|r|}{2u}$；$\beta'=d\nu-1$、$\gamma'=\alpha'=2-d\nu$、$d_c=4$ & W1005\\
\bottomrule
\end{tabular}
\end{center}
\textbf{筆誤一覽}（給 2004 年版；〔注〕見各節）：下表在文中都用紅框標出，每一條都用程式核對過。
\begin{center}
\footnotesize\renewcommand{\arraystretch}{0.9}
\TenTabErrata
\end{center}
